% problem-000077 1 主族元素与分子对称性
\section*{1. 主族元素与分子对称性}
图1-1⽰出了⼀种结构有趣的“糖葫芦”分⼦。其中A、B、C和D四种元素位于同⼀族的相邻周期，B的电负性⼤于 $\smash { \mathbf { D } _ { \circ } \textbf { C } }$ 可分别与A、B、D形成⼆元化合物，C和E可形成易挥发的化合物 $\mathrm { { C E } _ { 3 } , \ C E _ { 3 } }$ 在⼯业上可⽤于漂⽩和杀菌。A和E可形成化合物 $\mathbf { A E } _ { 3 } \left( \mathrm { b p , 7 6 ^ { \circ } C } \right)$ $\mathbf { A E } _ { 3 }$ 的完全⽔解产物之⼀X是塑料件镀⾦属的常⽤还原剂。A、C和E可形成含有六元坏的化合物 $\mathbf { A } _ { 3 } \mathbf { C } _ { 3 } \mathbf { E } _ { 6 \circ }$ 在氯仿中，D和E可结合两个甲基形成化合物 $\mathrm { { \bf D E } _ { 3 } M e _ { 2 } }$

（图：）
图 1-1

\subsection*{1-1}
写出A、B、C、D和E代表的元素。

\subsection*{1-2}
写出“糖葫芦”分⼦中A和C的杂化轨道。

\subsection*{1-3}
画出处于最⾼氧化态的A与元素E形成的分⼦结构图，写出该分⼦的对称元素（具体类型）。

\subsection*{1-4}
画出 $\mathrm { { \bf D E } _ { 3 } M e _ { 2 } }$ 的所有可能的⽴体异构体。

\subsection*{1-5}
指出X属于⼏元酸，写出弱酸介质中硫酸镍被其还原的反应⽅程式。
